\documentclass{article}
\usepackage[utf8]{inputenc}
\usepackage{amsmath}
\usepackage{graphicx}
\usepackage{geometry}
\geometry{a4paper, margin=1in}

\title{Homework 6 - Part 1a: Incompressible Taper-Flat Slider}
\author{}
\date{}

\begin{document}
\maketitle

\section*{Results}
For Part 1a, the pressure distribution for the 2D incompressible taper-flat slider bearing was computed using the Finite Element Method. 
The load $F$ and the center-of-pressure $x_p$ were calculated using the following integrals over the domain $\Omega$:
\begin{align*}
    F &= \int_{\Omega} p(x,y) \, dx \, dy \\
    x_p &= \frac{1}{F} \int_{\Omega} x \, p(x,y) \, dx \, dy
\end{align*}

Using the numerical solution, the computed values are:
\begin{itemize}
    \item \textbf{Computed Load ($F$):} $0.2066599$ N
    \item \textbf{Center of Pressure ($x_p$):} $0.002208667$ m
\end{itemize}

\section*{Analytical Verification}
The analytical load $\tilde{F}$ for a 1D pure wedge configuration is given by:
\begin{equation*}
    \tilde{F} = \frac{\Lambda L W}{(1-\alpha)^2} \left( \ln(\alpha) + 2 \frac{1-\alpha}{1+\alpha} \right)
\end{equation*}
where $\alpha = h_1/h_2$ and $\Lambda = 6\mu U L / h_2^2$.

Substituting the given parameters, we obtain:
\begin{itemize}
    \item \textbf{Analytical Load ($\tilde{F}$):} $8.154767$ N
\end{itemize}

\vspace{0.5cm}
\noindent\textit{Note:} The significant difference between the numerical load ($0.207$ N) and the analytical load ($8.15$ N) is expected. The numerical model simulates a 2D slider with a $1^{\circ}$ leading taper, whereas the analytical formula assumes a 1D, perfectly flat, full-length wedge.

\section*{Pressure Distribution}
The numerical pressure distribution $p(x,y)$ over the slider domain is shown in Figure \ref{fig:pressure}.

\begin{figure}[ht!]
    \centering
    % Since your image is uploaded as part_1a.jpg, this will link it automatically
    \includegraphics[width=0.8\textwidth]{part_1a}
    \caption{Normalized pressure distribution for the incompressible taper-flat slider bearing.}
    \label{fig:pressure}
\end{figure}

\end{document}